% paper3_rotation_curves_outline.tex
% Flat Rotation Curves from Quantum Relative Entropy:
% Running Newton's Constant as Information Processing
% Author: Sheng-Kai Huang
% Status: DETAILED OUTLINE
% Compile with: pdflatex paper3_rotation_curves_outline.tex

\documentclass[prl,twocolumn,superscriptaddress,longbibliography,floatfix]{revtex4-2}

\usepackage{amsmath}
\usepackage{amssymb}
\usepackage{amsfonts}
\usepackage{amsthm}
\usepackage{bm}
\usepackage{hyperref}
\usepackage{booktabs}
\usepackage{graphicx}
\usepackage{xcolor}

\newtheorem{theorem}{Theorem}

% Convenience macros (same as Paper 1)
\newcommand{\kB}{k_{\mathrm{B}}}
\newcommand{\Tr}{\mathrm{Tr}}
\newcommand{\id}{\mathrm{id}}
\newcommand{\Petz}{\mathcal{R}}
\newcommand{\chan}{\mathcal{N}}
\newcommand{\hilb}{\mathcal{H}}
\newcommand{\code}{\mathcal{C}}
\newcommand{\KL}{D}
\newcommand{\ket}[1]{|#1\rangle}
\newcommand{\bra}[1]{\langle#1|}
\newcommand{\braket}[2]{\langle#1|#2\rangle}
\newcommand{\op}[2]{|#1\rangle\langle#2|}
\newcommand{\abs}[1]{\left|#1\right|}
\newcommand{\norm}[1]{\left\|#1\right\|}
\newcommand{\tauasym}{\tau}

% Outline-specific formatting
\newcommand{\known}{\textsc{[known]}}
\newcommand{\newsyn}{\textsc{[new synthesis]}}
\newcommand{\newform}{\textsc{[new formula]}}
\newcommand{\spec}{\textsc{[speculative]}}
\newcommand{\towrite}[1]{{\color{blue}\textbf{[TO WRITE:} #1\textbf{]}}}
\newcommand{\toderive}[1]{{\color{red}\textbf{[TO DERIVE:} #1\textbf{]}}}
\newcommand{\assumed}[1]{{\color{orange}\textbf{[ASSUMED:} #1\textbf{]}}}

\begin{document}

\title{Flat Rotation Curves from Quantum Relative Entropy:\\
Running Newton's Constant as Information Processing}

\author{Sheng-Kai Huang}
\email{akai@fawstudio.com}
\affiliation{Independent Researcher}

\date{\today}

\begin{abstract}
\towrite{Final abstract after calculations are complete.}

We show that the quantum relative entropy (QRE) framework developed
in Papers~1 and~2---where entropy production
$\Sigma = D(\rho\|\sigma) - D(\chan(\rho)\|\chan(\sigma))$ quantifies
the arrow of time---naturally explains flat galaxy rotation curves
without dark matter particles. The key insight combines two
established results: (i)~Dorau and Much's derivation of
semiclassical Einstein equations from QRE at a fixed scale
(PRL~2025), and (ii)~Casini and Huerta's proof that
renormalization group (RG) flow is equivalent to data processing
inequality (DPI) monotonicity of QRE. Together, these imply that
the gravitational coupling $G$ inherits a scale dependence
$G(k)$ from the RG flow of the gravitational quantum channel.
With the marginal anomalous dimension $\eta = 1$
(Kumar~2025), this yields a logarithmic correction
$\Phi(r) = -G_N M/r + (2G_N M k_*/\pi)\ln(r/r_0)$ whose
universal coefficient $2/\pi$ produces flat rotation curves
at $r \gg r_c \approx 37$~kpc. We reinterpret the SPARC
galaxy fits of Gubitosi et~al.\ (2024) within the $\Sigma$
framework, where the running parameter $\bar{\nu}$ is identified
with the QRE drop per RG step, transforming it from a free
parameter into an information-theoretic quantity. Predictions
for BIG-SPARC ($\sim 4000$ galaxies) and honest limitations
(CMB, Bullet Cluster) are discussed.
\end{abstract}

\maketitle

%=======================================================================
\section{Introduction}
\label{sec:intro}
%=======================================================================

\subsection*{Content outline}

\known\ \emph{Papers~1--2 recap.}---Paper~1~\cite{Paper1}
established the temporal asymmetry parameter
\begin{equation}
\tauasym = 1 - F\!\big(\rho,\,
  \tilde{\Petz}_{\sigma,\chan}(\chan(\rho))\big),
\label{eq:tau_def}
\end{equation}
quantifying the arrow of time as Petz recovery failure, with the
universal bound $\tauasym \leq 1 - \exp(-\Sigma/2)$. Paper~2~\cite{Paper2}
showed that in the strong gravitational field of a static,
spherically symmetric mass, the entropy production takes the form
$\Sigma_{\mathrm{grav}} = r_s/r$, connecting the Petz recovery fidelity
to the gravitational redshift:
$F_{\mathrm{grav}} = \sqrt{g_{00}} = \exp(-r_s/(2r))$.

\known\ \emph{The dark matter puzzle.}---Galaxy rotation curves
flatten at large radii: $v(r) \to V_0 = \mathrm{const}$ for
$r \gtrsim 10$~kpc~\cite{Rubin1980,Bosma1981}, implying
$M_{\mathrm{dyn}}(r) \propto r$, far exceeding the visible baryonic
mass $M_{\mathrm{bar}}$. The standard explanation posits invisible
dark matter (DM) particles~\cite{Bertone2018}, but despite
four decades of searches, no DM particle has been detected
directly~\cite{Schumann2019,XENON2023}.

\newsyn\ \emph{Our proposal.}---We argue that flat rotation curves
emerge naturally from the \emph{same} $\Sigma = D(\rho\|\sigma) -
D(\chan(\rho)\|\chan(\sigma))$ framework. The gravitational quantum
channel $\chan_{\mathrm{grav}}$ is not scale-independent: quantum
corrections make $\chan_{\mathrm{grav}}(k)$ ``noisier'' at low
momenta $k$ (large distances), causing $G$ to run. This is not
a new term added by hand, but the natural consequence of evaluating
the \emph{same} $\Sigma$ formula with the full quantum-corrected
gravitational channel.

\towrite{Paragraph on analogy to $F = ma$ with different force laws
giving different orbits, springs, fluids. One equation, different
boundary conditions.}

\towrite{Paragraph on paper structure: Sec.~II reviews Dorau-Much;
Sec.~III connects RG to DPI; Sec.~IV derives running $G$ and
logarithmic correction; Sec.~V compares with SPARC data;
Sec.~VI gives predictions and limitations; Sec.~VII discusses
implications.}

%=======================================================================
\section{From Quantum Relative Entropy to Einstein Equations}
\label{sec:dorau_much}
%=======================================================================

\subsection*{Content outline}

\known\ \emph{Araki--Uhlmann modular theory.}---For a von~Neumann
algebra $\mathcal{M}$ with two faithful normal states $\omega_0$
(vacuum) and $\omega_\phi$ (coherent excitation), the Araki relative
entropy is~\cite{Araki1976}
\begin{equation}
S(\omega_\phi \| \omega_0)
= -\langle \Omega_\phi, \ln \Delta_{\omega_\phi|\omega_0}\,
  \Omega_\phi \rangle,
\label{eq:araki_rel_ent}
\end{equation}
where $\Delta_{\omega_\phi|\omega_0}$ is the relative modular
operator. This is the algebraic QFT generalization of the
finite-dimensional quantum relative entropy
$D(\rho\|\sigma) = \Tr[\rho(\ln\rho - \ln\sigma)]$.

\known\ \emph{Dorau--Much derivation (PRL~2025).}---The derivation
proceeds in five steps~\cite{DorauMuch2025}:
\begin{enumerate}
\item Start with a local Rindler horizon (equivalence principle):
  any point in spacetime admits an approximate Rindler wedge.
\item Compute $S(\omega_\phi \| \omega_0)$ between the vacuum
  $\omega_0$ and a coherent excitation $\omega_\phi$ of a scalar
  field on the right Rindler wedge:
  \begin{equation}
  S(\omega_\phi \| \omega_0)
  = -2\pi \int_{\mathcal{H}_B^R} U\,(\partial_U \phi)^2\, dU\, d\mathrm{vol}_S.
  \label{eq:dorau_much_srel}
  \end{equation}
\item Identify the area variation via the Bekenstein--Hawking
  relation: $\delta A = 4G\, S(\omega_\phi \| \omega_0)$.
\item Apply the Raychaudhuri equation to relate $\delta A$ to
  the Ricci tensor.
\item \emph{Result}: the semiclassical Einstein equations follow:
  \begin{equation}
  \boxed{
  G_{ab} + \Lambda g_{ab} = 8\pi G\, \langle T_{ab} \rangle.
  }
  \label{eq:einstein_from_qre}
  \end{equation}
\end{enumerate}

\towrite{Paragraph on significance: This establishes that QRE is
the \emph{microscopic origin} of Einstein's equations, extending
Jacobson's 1995 thermodynamic derivation~\cite{Jacobson1995} to
the quantum information setting. Crucially, this derivation holds
at a \emph{fixed scale}---it does not account for quantum
corrections that appear when the theory is probed at different
energy/distance scales.}

\emph{The fixed-scale limitation.}---Eq.~\eqref{eq:einstein_from_qre}
uses the \emph{bare} Newton constant $G$. It does not incorporate
the scale dependence of $G$ that arises from quantum loop corrections
to the gravitational channel. Our central claim is that this scale
dependence, when included via the data processing inequality,
naturally produces running $G(k)$ and flat rotation curves.

%=======================================================================
\section{Renormalization Group Flow as Data Processing}
\label{sec:rg_dpi}
%=======================================================================

\subsection*{Content outline}

\known\ \emph{Data processing inequality (DPI).}---For any
quantum channel $\chan$ and states $\rho, \sigma$~\cite{Lindblad1975,Uhlmann1977},
\begin{equation}
D\!\big(\chan(\rho)\big\|\chan(\sigma)\big)
\leq D(\rho\|\sigma).
\label{eq:dpi}
\end{equation}
Equality holds if and only if the Petz recovery map
$\Petz_{\sigma,\chan}$ exactly recovers $\rho$~\cite{Petz1986,Petz1988}.

\known\ \emph{RG flow as a quantum channel.}---Casini and
Huerta~\cite{CasiniHuerta2017} proved that renormalization group
(RG) flow from UV to IR can be formulated as a quantum channel:
the operation of integrating out high-energy modes is a CPTP map
$\chan_{\mathrm{RG}}: \rho_{\mathrm{UV}} \mapsto \rho_{\mathrm{IR}}$.
By the DPI,
\begin{equation}
D(\rho_{\mathrm{IR}} \| \sigma_{\mathrm{IR}})
\leq D(\rho_{\mathrm{UV}} \| \sigma_{\mathrm{UV}}).
\label{eq:rg_dpi}
\end{equation}
This monotonicity of QRE under RG flow provides:
\begin{itemize}
\item An alternative proof of the $c$-theorem in $d = 2$
  (Zamolodchikov);
\item In $d > 2$: the coefficient of the area-law term in
  entanglement entropy \emph{decreases} along RG flow;
\item A unified ``entropic'' proof of irreversibility theorems
  ($a$-theorem in $d = 4$, $F$-theorem in $d = 3$).
\end{itemize}

\newsyn\ \emph{Scale-dependent $\Sigma$.}---Define the entropy
production at scale $k$:
\begin{equation}
\Sigma(k) = D(\rho\|\sigma)
  - D\!\big(\chan_{\mathrm{grav}}(k)(\rho)
    \big\|\chan_{\mathrm{grav}}(k)(\sigma)\big).
\label{eq:sigma_k}
\end{equation}
Here $\chan_{\mathrm{grav}}(k)$ is the gravitational channel
evaluated at RG scale $k$: it describes the information processing
of gravity \emph{as seen at resolution $1/k$}.

\newsyn\ \emph{Key observation: as $k$ decreases (probing larger
distances), $\chan_{\mathrm{grav}}(k)$ becomes a ``noisier'' channel.}
This is because integrating out more UV modes increases the effective
dissipation. By the DPI, $\Sigma(k)$ can only increase or stay
the same as $k$ decreases:
\begin{equation}
\Sigma(k_1) \leq \Sigma(k_2) \quad \text{for } k_1 > k_2.
\label{eq:sigma_monotone}
\end{equation}
More entropy production at larger scales = more information loss =
stronger effective gravity. This is the information-theoretic origin
of running $G$.

\towrite{Paragraph connecting to Casini-Huerta's $c$-theorem:
the decrease of QRE along RG flow is the \emph{same} mathematical
structure as the DPI for the gravitational channel. RG irreversibility
\emph{is} data processing.}

\newsyn\ \emph{Observer dependence.}---Basso, Maziero, and
Celeri~\cite{BassoCeleri2025} proved that entropy production
in curved spacetime is observer-dependent, validating the
coordinate dependence of $\Sigma(k)$. De~Vuyst
et~al.~\cite{DeVuyst2025} confirmed this via quantum reference
frames. This is not a weakness of the framework but a
\emph{physical feature}: different observers probe gravity
at different scales and see different effective channels
$\chan_{\mathrm{grav}}(k)$.

Galley, Giacomini, and Selby~\cite{GalleyGiacominiSelby2023}
proved that \emph{any} consistent coupling between classical
gravity and quantum matter is fundamentally irreversible
($\Sigma > 0$). This provides model-independent justification
for the premise that gravitational channels produce entropy,
independent of the specific form of $\chan_{\mathrm{grav}}$.

%=======================================================================
\section{Running $G$ and the Logarithmic Correction}
\label{sec:running_G}
%=======================================================================

\subsection*{Content outline}

\known\ \emph{IR running of Newton's constant.}---In the effective
field theory of gravity, quantum loop corrections induce a
scale-dependent Newton coupling~\cite{Reuter1998,ReuterWeyer2004}:
\begin{equation}
G(k) \sim G_N \!\left(\frac{k_*}{k}\right)^{\!\eta}
\quad \text{for } k \ll k_*,
\label{eq:running_G}
\end{equation}
where $\eta$ is the anomalous dimension governing the IR flow and
$k_*$ is a crossover momentum scale.

\known\ \emph{Why $\eta = 1$ is special (marginal).}---Kumar~\cite{Kumar2025}
showed that the Fourier transform of $k^{-(2+\eta)}$ in three
spatial dimensions gives:
\begin{itemize}
\item $\eta \neq 1$: power-law $\sim r^{\eta - 1}$ (IR-irrelevant
  for $\eta < 1$, spoils scale invariance for $\eta > 1$);
\item $\eta = 1$: logarithmic $\sim \ln(r)$ --- the \emph{unique}
  marginal, scale-invariant deformation consistent with rotational
  symmetry and locality.
\end{itemize}

\assumed{$\eta = 1$ from marginality/dimensional analysis. Not yet
derived from QRE channel properties. This is the main theoretical
gap. See Sec.~\ref{sec:gaps} for discussion.}

\known\ \emph{Modified Newtonian potential.}---With
$G(k) = G_N[1 + (k_*/k)]$ for $\eta = 1$, the gravitational
potential in position space is~\cite{Kumar2025}
\begin{equation}
\boxed{
\Phi(r) = -\frac{G_N M}{r}
  + \frac{2G_N M k_*}{\pi}\,\ln\!\left(\frac{r}{r_0}\right),
}
\label{eq:phi_log}
\end{equation}
with $r_0$ an integration constant. The force law becomes
\begin{equation}
F(r) = -\frac{G_N M}{r^2}
  - \frac{2G_N M k_*}{\pi}\,\frac{1}{r}.
\label{eq:force_log}
\end{equation}
At $r \gg r_c = 1/k_*$, the $1/r$ force dominates over the
Newtonian $1/r^2$.

\known\ \emph{Flat rotation curves.}---The circular velocity is
\begin{equation}
v^2(r) = \frac{G_N M_{\mathrm{bar}}}{r}
  + \frac{2G_N M_{\mathrm{bar}} k_*}{\pi},
\label{eq:v_flat}
\end{equation}
which asymptotes to $v^2 \to 2G_N M_{\mathrm{bar}} k_*/\pi$
at $r \gg r_c$. The asymptotic velocity $V_0$ determines $k_*$
via~\cite{Kumar2025}
\begin{equation}
k_* = \frac{\pi V_0^2}{2G_N M_{\mathrm{bar}}}.
\label{eq:kstar}
\end{equation}

\known\ \emph{Universality of crossover scale.}---For three SPARC
galaxies, Kumar~\cite{Kumar2025} found:
\begin{center}
\begin{tabular}{lccc}
\toprule
Galaxy & $M_{\mathrm{bar}}\;[M_\odot]$ & $V_0$ [km/s]
  & $r_c$ [kpc] \\
\midrule
S:700624 & $1.01 \times 10^{12}$ & $271.4 \pm 2.6$ & $37.5 \pm 0.7$ \\
S:700916 & $5.85 \times 10^{11}$ & $210.4 \pm 2.4$ & $36.2 \pm 0.8$ \\
S:705253 & $3.55 \times 10^{11}$ & $158.9 \pm 1.9$ & $38.5 \pm 0.9$ \\
\bottomrule
\end{tabular}
\end{center}
The remarkably consistent $r_c \approx 37$~kpc across galaxies
spanning a factor of 3 in mass suggests a universal IR scale.

\known\ \emph{Free parameters.}---Only 2 parameters at the
galactic scale: $k_*$ (crossover scale) and $G_N$ (Newton's
constant). The coefficient $2/\pi$ is universal and
regulator-independent.

\newsyn\ \emph{Information-theoretic interpretation.}---In the
$\Sigma$ language, Eq.~\eqref{eq:phi_log} translates to
\begin{equation}
\Sigma_{\mathrm{grav}}(r) = \frac{r_s}{r}
  + \frac{4k_* r_s}{\pi}\,\ln\!\left(\frac{r}{r_0}\right),
\label{eq:sigma_grav_log}
\end{equation}
where $r_s = 2G_N M/c^2$. The logarithmic term represents
\emph{additional information loss} due to quantum corrections
to the gravitational channel at large distances. The channel
$\chan_{\mathrm{grav}}(k)$ at low $k$ (large $r$) loses more
information per unit QRE, producing extra $\Sigma$, which
manifests as extra gravitational potential.

\toderive{Show explicitly that $\Sigma(k)$ computed from the
Dorau--Much framework at each RG scale $k$ reproduces
Eq.~\eqref{eq:sigma_grav_log}. This requires defining
$\chan_{\mathrm{grav}}(k)$ using the RG-improved Hamiltonian
$H(k)$ with $G(k)$ and computing the QRE drop. This is the
central calculation of the paper.}

\toderive{Show that the coefficient $2/\pi$ emerges from the
channel properties of $\chan_{\mathrm{grav}}$ at the marginal
point $\eta = 1$. This would promote the coefficient from
``universal in Kumar's EFT sense'' to ``universal in the
information-theoretic sense.''}

%=======================================================================
\section{Comparison with SPARC Data}
\label{sec:sparc}
%=======================================================================

\subsection*{Content outline}

\known\ \emph{The SPARC database.}---The Spitzer Photometry
and Accurate Rotation Curves (SPARC) catalog~\cite{SPARC2016}
contains 175 nearby disk galaxies with 3.6~$\mu$m photometry
and high-quality rotation curves, providing the standard
benchmark for testing gravity theories at galactic scales.

\known\ \emph{RGGR fits.}---Gubitosi, Piattella, and
Casarini~\cite{Gubitosi2024} fitted 100 SPARC galaxies using
the renormalization group improved gravity (RGGR) model with
running coupling
\begin{equation}
G(\mu) = \frac{G_0}{1 + \nu\,\ln(\mu^2/\mu_0^2)},
\label{eq:G_RGGR}
\end{equation}
giving effective circular velocity
\begin{equation}
v_{\mathrm{RGGR}}^2(r) = v_N^2(r)
  \!\left[1 - \frac{c^2 \bar{\nu}}{\Phi_N(r)}\right],
\label{eq:v_RGGR}
\end{equation}
where $\bar{\nu} = \nu\alpha$ is a phenomenological parameter.
Using BIC analysis: 47 galaxies prefer RGGR over NFW,
40 prefer NFW, 13 inconclusive.

\known\ \emph{Mass-dependent $\bar{\nu}$.}---The RGGR parameter
$\bar{\nu}$ correlates near-linearly with galactic baryonic mass:
\begin{center}
\begin{tabular}{lc}
\toprule
Galaxy type & $\bar{\nu}$ \\
\midrule
Early type & $\sim 10^{-6}$ \\
Spiral & $\sim 10^{-7}$ \\
Late/Starburst & $\sim 10^{-8}$ \\
Solar system & $\sim 10^{-17}$ \\
\bottomrule
\end{tabular}
\end{center}

\newsyn\ \emph{Reinterpretation within $\Sigma$ framework.}---The
RGGR parameter $\bar{\nu}$ is identified with the QRE drop per
RG step of the gravitational channel:
\begin{equation}
\bar{\nu} = \frac{\Delta\Sigma}{\Delta\ln k^2}
  \bigg|_{k \sim k_*}\,,
\label{eq:nu_sigma}
\end{equation}
transforming $\bar{\nu}$ from a free parameter into an
information-theoretic observable. The mass dependence of
$\bar{\nu}$ then follows from the mass dependence of
$\Sigma_{\mathrm{grav}}$: more massive galaxies have larger
$r_s$, hence larger QRE drop per RG step.

\toderive{Compute $\Delta\Sigma/\Delta\ln k^2$ from the
gravitational channel and verify that the mass scaling matches
the observed $\bar{\nu}(M_{\mathrm{bar}})$ correlation.}

\known\ \emph{The Radial Acceleration Relation (RAR).}---The
tightest empirical constraint is the RAR~\cite{McGaugh2016}:
\begin{equation}
g_{\mathrm{obs}} = \frac{g_{\mathrm{bar}}}{1 - \exp\!\left(
  -\sqrt{g_{\mathrm{bar}}/a_0}\right)},
\label{eq:RAR}
\end{equation}
with $a_0 = 1.20 \times 10^{-10}$~m/s$^2$, observed scatter
0.13~dex over 2693 data points in 153 galaxies.

\newsyn\ \emph{RAR from running $\Sigma$.}---The logarithmic
correction in Eq.~\eqref{eq:phi_log} produces an additional
acceleration $g_{\mathrm{extra}} = 2G_N M_{\mathrm{bar}}k_*/(\pi r)$,
which becomes comparable to $g_{\mathrm{bar}} = G_N M_{\mathrm{bar}}/r^2$
precisely when $g_{\mathrm{bar}} \lesssim 2k_* c^2/\pi$.
Identifying
\begin{equation}
a_0 \sim \frac{2k_* c^2}{\pi},
\label{eq:a0_kstar}
\end{equation}
we recover the correct order of magnitude for the MOND
acceleration scale. With $k_* \approx 2.7 \times 10^{-2}$~kpc$^{-1}$,
this gives $a_0 \sim 1.7 \times 10^{-10}$~m/s$^2$, within a
factor of 1.4 of the observed value.

\towrite{Paragraph on the coincidence $a_0 \sim cH_0$
(Milgrom~\cite{Milgrom2020}) and its possible information-theoretic
origin: the de Sitter horizon entropy sets the IR cutoff for the
gravitational channel, connecting the acceleration scale to
cosmology.}

%=======================================================================
\section{Predictions, Tests, and Limitations}
\label{sec:predictions}
%=======================================================================

\subsection*{Content outline}

\newsyn\ \emph{Prediction 1: BIG-SPARC.}---The upcoming BIG-SPARC
catalog~\cite{BIGSPARC2024} with $\sim 4000$ galaxies (20$\times$
larger than SPARC) provides a stringent test. Our prediction:
the crossover scale $r_c \approx 37$~kpc should remain universal
across the expanded sample, and the RGGR parameter $\bar{\nu}$
should maintain its near-linear correlation with baryonic mass.

\newsyn\ \emph{Prediction 2: Dwarf galaxies.}---Dwarf spheroidal
galaxies are the strongest test case, as they have the largest
mass discrepancies. Ghari and Haghi~\cite{GhariHaghi2026} found
that Verlinde's emergent gravity (which shares the same
information-theoretic structure as our approach) is favored over
MOND at $5.2\sigma$ for 23 dwarf spheroidals. Our prediction:
the $\Sigma$ framework should reproduce these fits with the
\emph{same} universal parameters.

\newsyn\ \emph{Prediction 3: Solar system screening.}---At
solar system scales ($r \sim 1$~AU $\ll r_c \approx 37$~kpc), the
logarithmic correction in Eq.~\eqref{eq:phi_log} contributes
$\delta\Phi/\Phi_N \sim 2k_* r/\pi \sim 10^{-15}$, far below
current PPN constraints. The theory is automatically screened at
small scales.

\known\ \emph{Limitation 1: CMB acoustic peaks.}---Running $G$
alone \emph{cannot} reproduce the CMB acoustic peak structure
without dark matter~\cite{Kumar2025}. The IR running effect is
negligible at $z \sim 1100$ when the universe was $\sim 380{,}000$
years old. This is a severe limitation.

\towrite{Honest discussion: Does this rule out the approach? Not
necessarily---the CMB constraint may require \emph{additional}
physics beyond running $G$ (e.g., sterile neutrinos, primordial
fluctuations, or modifications to the gravitational channel at
earlier epochs). The $\Sigma$ framework is open to incorporating
such effects if they arise naturally from the channel properties.
But we cannot claim to replace dark matter entirely.}

\known\ \emph{Limitation 2: Bullet Cluster.}---The Bullet Cluster
(1E~0657-56) shows spatial separation between the gravitational
potential (traced by lensing) and the baryonic gas (traced by
X-rays)~\cite{Clowe2006}. No systematic analysis of the $\Sigma$
framework at cluster scales has been performed. This is a
well-known challenge for all modified gravity theories.

\known\ \emph{Limitation 3: Structure formation.}---Cosmological
structure formation at high redshift requires matter perturbations
that grow faster than baryons alone can provide. Running $G$ may
help~\cite{ReuterWeyer2004}, but a full perturbation theory within
the $\Sigma$ framework has not been developed.

\towrite{Summary table of what the approach explains, what it
does not, and what remains open:}

\begin{table}[htbp]
\caption{\label{tab:status}Status of observational tests.}
\begin{ruledtabular}
\begin{tabular}{lcc}
\textbf{Observable} & \textbf{Status} & \textbf{Notes} \\
\colrule
Flat rotation curves & Yes & 3 galaxies (Kumar) \\
SPARC (100 galaxies) & Yes & RGGR fits (Gubitosi+) \\
RAR ($a_0$ scale) & Partial & Order-of-magnitude \\
Solar system PPN & Yes & Auto-screened \\
Dwarf spheroidals & Promising & Verlinde 5.2$\sigma$ \\
Weak lensing (KiDS) & Promising & Brouwer+ 2021 \\
\colrule
CMB acoustic peaks & \textbf{No} & Requires extra physics \\
Bullet Cluster & Unknown & Not analyzed \\
Structure formation & Unknown & No perturbation theory \\
Lyman-$\alpha$ forest & Unknown & Not analyzed \\
\end{tabular}
\end{ruledtabular}
\end{table}

%=======================================================================
\section{Key Theoretical Gaps}
\label{sec:gaps}
%=======================================================================

\subsection*{Content outline}

\towrite{Honest assessment of the three main gaps.}

\emph{Gap 1: Deriving $\eta = 1$ from QRE.}---The anomalous
dimension $\eta = 1$ is currently argued from dimensional analysis
and marginality~\cite{Kumar2025}, not derived from the channel
properties of $\chan_{\mathrm{grav}}$. A first-principles
derivation from QRE would require showing that $\eta = 1$ is the
unique value at which the DPI bound for the gravitational channel
is ``marginally'' saturated---i.e., the boundary between
recoverable and non-recoverable regimes of the channel.

\spec\ \emph{Possible argument}: In the information-theoretic
framework, $\eta < 1$ corresponds to a channel where the DPI
deficit \emph{decreases} with scale (the channel becomes ``less
noisy'' in the IR), which would mean gravity \emph{weakens} at
large distances---unphysical. $\eta > 1$ corresponds to
\emph{super-logarithmic} growth of $\Sigma$, spoiling the
perturbative expansion. $\eta = 1$ is the unique marginal
case where $\Sigma$ grows logarithmically, maintaining the
delicate balance between gravity strengthening (as required by
observations) and perturbative control.

\spec\ \emph{New conjecture}: Applying the Fawzi--Renner
recovery bound to the gravitational RG flow suggests
\begin{equation}
\eta \leq 1 - \exp(-\Sigma_{\mathrm{RG}}/2),
\label{eq:eta_bound}
\end{equation}
where $\Sigma_{\mathrm{RG}}$ is the entropy production per RG
step. For large $\Sigma_{\mathrm{RG}}$, this gives $\eta \leq 1$,
making $\eta = 1$ the maximal allowed value. In flat space
($\Sigma_{\mathrm{RG}} = 0$), $\eta = 0$, which is the correct
UV fixed point.

\emph{Gap 2: Defining $\chan_{\mathrm{grav}}(k)$ rigorously.}---The
gravitational channel at scale $k$ can be defined via:
\begin{itemize}
\item Option A: RG-improved Hamiltonian $H(k)$ with $G(k)$
  replacing $G_N$ in the Pikovski-type coupling;
\item Option B: One-loop corrected graviton propagator in the
  effective field theory;
\item Option C: Modular channel with scale-dependent modular
  Hamiltonian (most rigorous but hardest to compute).
\end{itemize}
A rigorous definition is needed to compute $\Sigma(k)$ explicitly.

\emph{Gap 3: The acceleration scale $a_0 \sim cH_0$.}---The
coincidence between the MOND acceleration scale
$a_0 \approx 1.2 \times 10^{-10}$~m/s$^2$ and $cH_0 \approx
6.9 \times 10^{-10}$~m/s$^2$ is suggestive but unexplained
within the $\Sigma$ framework. Possible origin: the de~Sitter
horizon entropy $S_{\mathrm{dS}} = \pi c^3/(G\hbar H_0^2)$ sets
the IR cutoff for the gravitational channel, making
$k_{\mathrm{min}} \sim H_0/c$ and $a_0 \sim c^2 k_{\mathrm{min}}$.
This connects galactic dynamics to cosmology through the
\emph{same} QRE structure.

%=======================================================================
\section{Discussion}
\label{sec:discussion}
%=======================================================================

\subsection*{Content outline}

\newsyn\ \emph{What is genuinely new.}---Three things:
\begin{enumerate}
\item \emph{The connection}: QRE at each scale $\to$ running $G$
  $\to$ rotation curves. Dorau-Much~\cite{DorauMuch2025} gives
  Einstein equations from QRE at a fixed scale. Casini-Huerta~\cite{CasiniHuerta2017}
  shows QRE monotonicity = RG irreversibility. We combine them:
  QRE at each scale gives scale-dependent Einstein equations
  with running $G$.
\item \emph{The reinterpretation}: The RGGR parameter $\bar{\nu}$
  is not a free parameter but the QRE drop per RG step, connecting
  phenomenological galaxy fitting to fundamental information theory.
\item \emph{The unification}: Papers~1, 2, and 3 use the \emph{same}
  $\Sigma$ framework at different scales. $\Sigma = 0$ (Paper~1,
  QI), $\Sigma = r_s/r$ (Paper~2, strong field),
  $\Sigma = r_s/r + \text{log correction}$ (Paper~3, weak field).
\end{enumerate}

\emph{What is NOT new.}---Running $G$ has been proposed since
Reuter and Weyer (2004)~\cite{ReuterWeyer2004}. Kumar~\cite{Kumar2025}
derived the logarithmic correction. Gubitosi et~al.~\cite{Gubitosi2024}
fitted SPARC galaxies. Our contribution is the
\emph{information-theoretic grounding}: why $G$ runs, what the
running means in terms of quantum channels, and how it connects
to the broader $\Sigma$ framework.

\towrite{Paragraph on what this does NOT achieve: we do not explain
CMB, Bullet Cluster, or structure formation. We are honest that
``replacing dark matter'' is too strong a claim. What we show is
that \emph{at galactic scales}, running $\Sigma$ from QRE provides
a viable alternative that is competitive with NFW halos and has
a deeper theoretical motivation.}

\emph{Comparison with other information-theoretic approaches.}---
Several groups have connected information theory to dark matter:
\begin{itemize}
\item Verlinde~\cite{Verlinde2016}: Volume-law entanglement entropy
  in de~Sitter $\to$ apparent dark matter. Complementary to our
  approach; may be the same physics in different language.
\item Bianconi~\cite{Bianconi2025}: Gravity action = QRE,
  metric as density matrix. However, the GfE vacuum solutions
  give Schwarzschild + $O(r^{-4})$ corrections~\cite{Thattarampilly2026},
  not the logarithmic modification needed at galactic scales.
  Our approach via running $\Sigma(k)$ does not depend on GfE.
\item Smolin~\cite{Smolin2017}: MOND from quantum gravity regime
  below de~Sitter temperature.
\end{itemize}

\newsyn\ \emph{Connection to Papers~1--2 and preview of Paper~4.}---
The logical chain is now:
\begin{align}
&\text{Paper 1:}\;\; \Sigma = 0 \Leftrightarrow \text{perfect recovery},
  \nonumber\\
&\text{Paper 2:}\;\; \Sigma_{\mathrm{grav}} = r_s/r
  \;\to\; \text{exponential metric},
  \nonumber\\
&\text{Paper 3:}\;\; \Sigma_{\mathrm{grav}}(k) = r_s(k)/r
  \;\to\; \text{flat curves},
  \nonumber\\
&\text{Paper 4:}\;\; \Sigma = D(\rho_{\mathrm{st}} \| \rho_{\mathrm{m}})
  \;\to\; \text{all three as limits}.
  \nonumber
\end{align}
Paper~4~\cite{Paper4} will argue that these are not three
separate theories but three solutions of \emph{one equation}
$\Sigma = D(\rho_{\mathrm{spacetime}} \| \rho_{\mathrm{matter}})$
under different boundary conditions, analogous to $F = ma$ giving
planetary orbits, spring oscillations, and fluid dynamics.

\begin{acknowledgments}
The author thanks the quantum information, quantum gravity,
and astrophysics communities for the foundational works on
which this synthesis builds. Special thanks to the SPARC team
for making their data publicly available.
\end{acknowledgments}

%-----------------------------------------------------------------------
% BIBLIOGRAPHY
%-----------------------------------------------------------------------

\begin{thebibliography}{50}

\bibitem{Paper1}
S.-K. Huang,
``Universal recovery: The Petz map as retrodiction functor,
error corrector, and entropy bound,''
companion paper (2026).

\bibitem{Paper2}
S.-K. Huang,
``Information-theoretic origin of the exponential metric:
Gravitational redshift as Petz recovery fidelity,''
companion paper (2026).

\bibitem{Paper4}
S.-K. Huang,
``One equation, three scales: $\Sigma = D(\rho_{\mathrm{spacetime}}
\| \rho_{\mathrm{matter}})$ as the universal action,''
in preparation (2026).

\bibitem{DorauMuch2025}
M.~Dorau and A.~Much,
``From quantum relative entropy to the semiclassical Einstein
equations,''
Phys.\ Rev.\ Lett.\ (2025); arXiv:2510.24491.

\bibitem{CasiniHuerta2017}
H.~Casini and M.~Huerta,
``Relative entropy and the RG flow,''
J.\ High Energy Phys.\ \textbf{2017}, 044 (2017);
arXiv:1611.00016.

\bibitem{Kumar2025}
N.~Kumar,
``Marginal IR running of gravity as a natural explanation for
dark matter,''
Phys.\ Lett.\ B \textbf{871}, 140008 (2025);
arXiv:2509.05246.

\bibitem{Gubitosi2024}
G.~Gubitosi, O.~F. Piattella, and R.~M. Casarini,
``Phenomenology of renormalization group improved gravity from
the kinematics of SPARC galaxies,''
(2024); arXiv:2403.00531.

\bibitem{Rubin1980}
V.~C. Rubin, W.~K. Ford, Jr., and N.~Thonnard,
``Rotational properties of 21 SC galaxies with a large range
of luminosities and radii,''
Astrophys.\ J.\ \textbf{238}, 471 (1980).

\bibitem{Bosma1981}
A.~Bosma,
``21-cm line studies of spiral galaxies. II: The distribution
and kinematics of neutral hydrogen,''
Astron.\ J.\ \textbf{86}, 1825 (1981).

\bibitem{Bertone2018}
G.~Bertone and D.~Hooper,
``History of dark matter,''
Rev.\ Mod.\ Phys.\ \textbf{90}, 045002 (2018).

\bibitem{Schumann2019}
M.~Schumann,
``Direct detection of WIMP dark matter: Concepts and status,''
J.\ Phys.\ G \textbf{46}, 103003 (2019).

\bibitem{XENON2023}
XENON Collaboration,
``First dark matter search with nuclear recoils from the
XENONnT experiment,''
Phys.\ Rev.\ Lett.\ \textbf{131}, 041003 (2023).

\bibitem{Petz1986}
D.~Petz,
``Sufficient subalgebras and the relative entropy of states of a
von Neumann algebra,''
Commun.\ Math.\ Phys.\ \textbf{105}, 123 (1986).

\bibitem{Petz1988}
D.~Petz,
``Sufficiency of channels over von Neumann algebras,''
Q.\ J.\ Math.\ \textbf{39}, 97 (1988).

\bibitem{Lindblad1975}
G.~Lindblad,
``Completely positive maps and entropy inequalities,''
Commun.\ Math.\ Phys.\ \textbf{40}, 147 (1975).

\bibitem{Uhlmann1977}
A.~Uhlmann,
``Relative entropy and the Wigner-Yanase-Dyson-Lieb concavity
in an interpolation theory,''
Commun.\ Math.\ Phys.\ \textbf{54}, 21 (1977).

\bibitem{Araki1976}
H.~Araki,
``Relative entropy of states of von Neumann algebras,''
Publ.\ RIMS Kyoto \textbf{11}, 809 (1976).

\bibitem{Jacobson1995}
T.~Jacobson,
``Thermodynamics of spacetime: The Einstein equation of state,''
Phys.\ Rev.\ Lett.\ \textbf{75}, 1260 (1995).

\bibitem{Reuter1998}
M.~Reuter,
``Nonperturbative evolution equation for quantum gravity,''
Phys.\ Rev.\ D \textbf{57}, 971 (1998).

\bibitem{ReuterWeyer2004}
M.~Reuter and H.~Weyer,
``Running Newton constant, improved gravitational actions, and
galaxy rotation curves,''
Phys.\ Rev.\ D \textbf{70}, 124028 (2004).

\bibitem{SPARC2016}
F.~Lelli, S.~S. McGaugh, and J.~M. Schombert,
``SPARC: Mass models for 175 disk galaxies with Spitzer
photometry and accurate rotation curves,''
Astron.\ J.\ \textbf{152}, 157 (2016); arXiv:1606.09251.

\bibitem{McGaugh2016}
S.~S. McGaugh, F.~Lelli, and J.~M. Schombert,
``Radial acceleration relation in rotationally supported
galaxies,''
Phys.\ Rev.\ Lett.\ \textbf{117}, 201101 (2016);
arXiv:1609.05917.

\bibitem{BIGSPARC2024}
L.~Haubner, F.~Lelli, \emph{et~al.},
``BIG-SPARC: BCG, isolated, and group Spitzer photometry and
accurate rotation curves,''
(2024); arXiv:2411.13329.

\bibitem{GhariHaghi2026}
A.~Ghari and H.~Haghi,
``Testing emergent gravity in dwarf spheroidal galaxies,''
(2026); arXiv:2601.01715.

\bibitem{Milgrom2020}
M.~Milgrom,
``MOND vs.\ dark matter in light of historical parallels,''
(2020); arXiv:2001.09729.

\bibitem{Clowe2006}
D.~Clowe, M.~Brada\v{c}, A.~H. Gonzalez, M.~Markevitch,
S.~W. Randall, C.~Jones, and D.~Zaritsky,
``A direct empirical proof of the existence of dark matter,''
Astrophys.\ J.\ Lett.\ \textbf{648}, L109 (2006).

\bibitem{Verlinde2016}
E.~P. Verlinde,
``Emergent gravity and the dark universe,''
SciPost Phys.\ \textbf{2}, 016 (2017); arXiv:1611.02269.

\bibitem{Bianconi2025}
G.~Bianconi,
``Gravity from entropy,''
Phys.\ Rev.\ D \textbf{111}, 066001 (2025);
arXiv:2408.14391.

\bibitem{Smolin2017}
L.~Smolin,
``MOND as a regime of quantum gravity,''
Phys.\ Rev.\ D \textbf{96}, 083523 (2017);
arXiv:1704.00780.

\bibitem{Rodrigues2009}
D.~C. Rodrigues, P.~S. Letelier, and I.~L. Shapiro,
``Galaxy rotation curves from general relativity with infrared
renormalization group effects,''
J.\ Cosmol.\ Astropart.\ Phys.\ \textbf{2010}, 020 (2010);
arXiv:0911.4967.

\bibitem{Donoghue2020}
J.~F. Donoghue,
``A critique of the asymptotic safety program,''
Front.\ Phys.\ \textbf{8}, 56 (2020).

\bibitem{Li2018}
P.~Li, F.~Lelli, S.~S. McGaugh, and J.~M. Schombert,
``Fitting the radial acceleration relation to individual SPARC
galaxies,''
Astron.\ Astrophys.\ \textbf{615}, A3 (2018);
arXiv:1803.00022.

\bibitem{BrouwerKiDS2021}
M.~M. Brouwer \emph{et~al.},
``The weak lensing radial acceleration relation,''
Astron.\ Astrophys.\ \textbf{650}, A113 (2021);
arXiv:2106.11677.

\bibitem{CasiniHuerta2004}
H.~Casini and M.~Huerta,
``A finite entanglement entropy and the $c$-theorem,''
Phys.\ Lett.\ B \textbf{600}, 142 (2004);
hep-th/0405111.

\bibitem{Jacobson2016}
T.~Jacobson,
``Entanglement equilibrium and the Einstein equation,''
Phys.\ Rev.\ Lett.\ \textbf{116}, 201101 (2016);
arXiv:1505.04753.

\bibitem{BassoCeleri2025}
M.~L.~W.~Basso, J.~Maziero, and L.~C.~Celeri,
``Quantum detailed fluctuation theorem in curved spacetime,''
Phys.\ Rev.\ Lett.\ \textbf{134}, 050406 (2025);
arXiv:2405.03902.

\bibitem{GalleyGiacominiSelby2023}
T.~D.~Galley, F.~Giacomini, and J.~O.~E.~Selby,
``Any consistent coupling between classical gravity and quantum matter is fundamentally irreversible,''
Quantum \textbf{7}, 1142 (2023).

\bibitem{DeVuyst2025}
J.~De~Vuyst, S.~Eccles, P.~A.~H\"{o}hn, and J.~Kirklin,
``Gravitational entropy is observer-dependent,''
(2025); arXiv:2412.15502.

\bibitem{Thattarampilly2026}
B.~Thattarampilly, Y.~Zheng, and K.~K.~Kakkat,
``Spherically symmetric black hole in gravity from entropy,''
(2026); arXiv:2602.13694.

\bibitem{KudlerFlam2024}
J.~Kudler-Flam,
``Fluctuation theorems from Connes-Takesaki (non-)commutative flow of weights,''
(2024); arXiv:2408.04219.

\end{thebibliography}

\end{document}
